\documentclass[11pt,a4paper]{article}
\usepackage[margin=1in]{geometry}
\usepackage{hyperref}
\usepackage{enumitem}
\usepackage{graphicx}
\usepackage{xcolor}

% -----------------------------------------------------
% bvraghav-tiet-ucs503p-article.sty
% -----------------------------------------------------

% Variable: Lab Instructor
\makeatletter \newcommand{\BvrSetLabInstructor}[1]{%
  \gdef\Bvr@LabInstructor{#1}}
% default
\newcommand{\Bvr@LabInstructor}{Dr. Raghav %
  B. Venkataramaiyer}
% public accessor
\newcommand{\BvrLabInstructorValue}{\Bvr@LabInstructor}
\makeatother

% Add subtitle
\usepackage{titling}
% -- Set title bold
\pretitle{\begin{center}\bfseries\fontsize{18bp}{18bp}\selectfont}
\makeatletter
\newcommand{\BvrSubtitle}[1]{%
  \posttitle{%
    \par\end{center}
    \begin{center}\large#1\end{center}
    \vskip 0.5em}%
}
\makeatother

% Hints in the section title(s)
\newcommand{\BvrHint}[1]{%
  {\normalfont\normalsize\itshape\hfill #1}}

% -----------------------------------------------------

\title{SV Car Customs: A Vehicle-Aware E-Commerce Platform for Car Accessories}

\BvrSetLabInstructor{Ms.\ Anushka Chaudhary}

\BvrSubtitle{%
  Project Proposal for UCS503P  \\
  Submitted to: \BvrLabInstructorValue \\ \bfseries
  Thapar Institute of Engineering and Technology%
}

\author{
  Vishu Verma, COPC\thanks{Roll No: \texttt{1024170242}}
  \and
  Raghu Verma, COPC\thanks{Roll No: \texttt{1024170241}}
  \and
  Pratham Sharma, COPC\thanks{Roll No: \texttt{1024170229}}
  \and
  Tarush Mohindru, COPC\thanks{Roll No: \texttt{1024170231}}
}

\date{Group \texttt{3Q24} \quad (revised 26 November 2026)}

\begin{document}
\maketitle

\section{Higher-order goal%
  \BvrHint{\ldots secondary framing}}
Help a small Indian retail business sell online on its own terms: own its
customer data, appear in Google search, and take payments safely, without
depending on a hosted no-code store. The immediate engineering objective is a
deployable e-commerce platform for SV Car Customs, a Delhi seller of
aftermarket car accessories.

\section{Time-to-value %
\BvrHint{\ldots primary heuristic}}
\begin{itemize}[leftmargin=0.7cm]
    \item The first increment is a browsable catalogue with product pages on a
          public URL, so the owner can share links within the first two weeks.
    \item Each following week adds one end-to-end slice (admin, login and cart,
          payments, search and fitment), shipped behind CI.
    \item Success in iteration 1 is measured by the owner being able to add a
          product in the admin panel and see it live on the storefront.
\end{itemize}

\section{Problem Statement}
SV Car Customs sells body kits, spoilers, diffusers, exhausts, lighting and
interior parts to car owners across India. Its existing hosted storefront
causes:
\begin{itemize}[leftmargin=0.7cm]
    \item \textbf{Wrong-fit orders:} products are not linked to the car makes
          and models they fit, so customers must ask on WhatsApp or guess.
    \item \textbf{Low search visibility:} little control over page titles,
          structured data or a Google Shopping feed.
    \item \textbf{Payment friction:} no Cash-on-Delivery option with an online
          advance, which Indian buyers of high-value parts expect.
    \item \textbf{Manual operations:} orders, stock and COD collection are
          tracked outside the store.
\end{itemize}

\textbf{Why this matters now:} the business already gets most of its traffic
from Instagram and wants to convert it into orders on its own site.

\section{Proposed Solution %
\BvrHint{\ldots Software Proposition}}
\subsection{Overview}
Build a \textbf{web-based} store with:
\begin{itemize}[leftmargin=0.7cm]
    \item \textbf{Storefront:} catalogue by car model and part type, product
          pages with variants and fitment, search, cart and checkout.
    \item \textbf{Vehicle fitment:} products linked to vehicle makes and models;
          a ``My Garage'' selector to show only parts that fit.
    \item \textbf{Payments:} PayU online payment, prepaid or COD with a 15\%
          online advance.
    \item \textbf{Admin panel:} products, categories, orders, COD collection,
          settings, home page and an audit log.
    \item \textbf{SEO:} sitemap, structured data and a Merchant Center feed.
\end{itemize}

\subsection{Core workflow}
\begin{enumerate}[leftmargin=0.7cm]
    \item Customer selects their car and browses parts that fit it.
    \item Customer adds a variant (SKU) to the cart and checks out.
    \item The server computes the total and creates the order, then redirects to
          PayU.
    \item The server verifies PayU's response and marks the order paid.
    \item The owner ships the order and updates its status; the customer tracks it.
\end{enumerate}

\subsection{Operational constraints}
\begin{itemize}[leftmargin=0.7cm]
    \item Mobile-first: most visitors arrive from Instagram on phones.
    \item Free or low-cost hosting tiers only.
    \item Product claims (fitment, material) must come from the owner's real
          data, never be invented.
\end{itemize}

\section{Solution Approach %
\BvrHint{\ldots Engineering focus}}
\subsection{Trusted pricing and payments %
\BvrHint{\ldots non-research}}
\begin{itemize}[leftmargin=0.7cm]
    \item The browser is untrusted: prices are recomputed on the server from the
          database for every checkout.
    \item Payment state changes only after the gateway's signed response is
          verified; a scheduled job reconciles orders whose callback was lost.
\end{itemize}

\subsection{Web architecture %
\BvrHint{\ldots web-based solution}}
\begin{itemize}[leftmargin=0.7cm]
    \item \textbf{Frontend:} Next.js (React, TypeScript) on Vercel, with
          server-rendered, cached pages.
    \item \textbf{Backend API:} Hono (Node.js) on Fly.io: authentication,
          catalogue, cart, checkout, payments and admin routes.
    \item \textbf{Database:} PostgreSQL and object storage on Supabase.
\end{itemize}

\subsection{CI/CD and fast delivery enabler}
\begin{itemize}[leftmargin=0.7cm]
    \item GitHub Actions runs type-checks, tests and builds for both apps on
          every pull request.
    \item The frontend deploys automatically from the main branch; the API is
          deployed as a Docker image.
    \item Direct pushes to the main branch are blocked; changes go through pull
          requests.
\end{itemize}

\section{Evaluation Criterion %
\BvrHint{\ldots measurable and attributable}}
\subsection{Primary evaluation metric}
\textbf{Checkout completion:} share of checkout views that end in a confirmed
order, measured from the store's own analytics events and order records.

\subsection{Secondary evaluation metrics}
\begin{itemize}[leftmargin=0.7cm]
    \item \textbf{Payment integrity:} zero paid orders left in \emph{pending}
          after reconciliation.
    \item \textbf{Speed:} Largest Contentful Paint under 2.5\,s on 4G; cached
          pages under 300\,ms.
    \item \textbf{Search visibility:} clicks and impressions in Google Search
          Console.
    \item \textbf{Quality:} all automated tests passing in CI on every merge.
\end{itemize}

\section{Scalability %
\BvrHint{\ldots with foresight}}
\begin{enumerate}[leftmargin=0.7cm]
    \item \textbf{Theoretically:} stateless API, cached pages, paginated
          category listings and indexed queries.
    \item \textbf{Practically:} managed database, container-hosted API and a
          white-label configuration file so the platform can be reused for
          another store.
\end{enumerate}

\section{Engine availability heuristic}
\begin{itemize}[leftmargin=0.7cm]
    \item Next.js and Hono for web pages and REST APIs.
    \item Supabase PostgreSQL and storage; Drizzle ORM.
    \item PayU for payments; Resend and MSG91 for email and SMS.
    \item Vitest and Mock Service Worker for testing; GitHub Actions for CI.
\end{itemize}

\section{Project scope and deliverables %
\BvrHint{\ldots iteration-friendly}}
\subsection{Initial deliverable %
\BvrHint{\ldots first iteration}}
\begin{itemize}[leftmargin=0.7cm]
    \item Database schema, catalogue API and storefront category/product pages
    \item Admin login and product management
    \item CI: type-check, tests and build for both apps
    \item Deployment of both apps
\end{itemize}

\subsection{Subsequent deliverables}
\begin{itemize}[leftmargin=0.7cm]
    \item Email-OTP login, cart and checkout
    \item PayU payments, COD advance and reconciliation job
    \item Vehicle fitment, search and ``My Garage''
    \item SEO, analytics and admin order management
\end{itemize}

\section{Risks and mitigations}
\begin{itemize}[leftmargin=0.7cm]
    \item \textbf{Payment callback loss:} webhook plus scheduled reconciliation.
    \item \textbf{Free-tier limits:} caching, image optimisation outside Vercel's
          quota, and monitoring of usage.
    \item \textbf{Inaccurate product data:} owner confirms fitment and materials
          before publishing.
    \item \textbf{Security:} hashed OTPs and passwords, rate limits, server-side
          validation, secrets only in environment variables.
\end{itemize}

\section{Summary}
The project replaces a hosted store with a deployable, vehicle-aware
e-commerce platform for a real business. It meets the course criteria through
weekly increments backed by CI, measurable evaluation (checkout completion,
payment integrity, speed and search visibility) and an architecture that can
scale and be reused.

\end{document}
